\documentclass[11pt,a4paper]{article}

\usepackage[T1]{fontenc}
\usepackage[utf8]{inputenc}
\usepackage{lmodern}
\usepackage{microtype}
\usepackage[margin=1in]{geometry}
\emergencystretch=3em
\usepackage{amsmath,amssymb,amsthm}
\usepackage{booktabs}
\usepackage{multirow}
\usepackage{array}
\usepackage{graphicx}
\usepackage{xcolor}
\usepackage{caption}
\usepackage{enumitem}
\usepackage{algorithm}
\usepackage{algpseudocode}
\usepackage{listings}
\usepackage[hidelinks,colorlinks=true,linkcolor=black,citecolor=black,urlcolor=blue]{hyperref}

\definecolor{codebg}{gray}{0.96}
\lstset{
  basicstyle=\ttfamily\small,
  backgroundcolor=\color{codebg},
  frame=none,
  breaklines=true,
  columns=fullflexible,
  keepspaces=true,
  showstringspaces=false,
  xleftmargin=1em,
  framexleftmargin=1em,
}

\newcommand{\code}[1]{\texttt{\small #1}}
\newcommand{\SXE}{\textsc{SXE}}
\DeclareMathOperator*{\argmax}{arg\,max}
\DeclareMathOperator*{\argmin}{arg\,min}
\DeclareMathOperator{\Var}{Var}
\DeclareMathOperator{\SSE}{SSE}

\theoremstyle{definition}
\newtheorem{observation}{Observation}

\title{\Large \SXE: Batch-Boosted, Holdout-Pruned, Sparsely Re-Weighted
Histogram Trees\\[0.4em]
\large A tree ensemble between the random forest and the gradient-boosted machine}

\author{Alastair Rushworth}

\date{%
  \today\\[0.5em]
  \small Prepared from the \code{sxe} repository at commit \code{eb97f58};
  every measured figure traces to a script or artefact named in the text.
}

\begin{document}
\maketitle

\begin{abstract}
\noindent
\SXE\ is a supervised tree-ensemble method for tabular data that occupies the
space between two familiar constructions. A random forest grows many
independent trees against the original target and averages them; a
gradient-boosted decision tree (GBDT) grows one tree at a time against a
running residual. \SXE\ grows a small \emph{batch} of independent, randomized
trees at each boosting step, averages the batch, and applies that average as a
single shrunken correction to the residual. The batch size is therefore a
continuous dial between the two classical methods: a batch of one recovers
stochastic gradient boosting, and a single batch spanning the whole budget
recovers a random forest.

Two further devices sit on top of that generator. First, because each tree is
grown on a row subsample, the complement of that subsample is a private
holdout, and \SXE\ uses it for reduced-error pruning at a strength drawn at
random \emph{per tree}, so the ensemble deliberately spans a spectrum of
bias--variance trade-offs rather than committing to one. Second, after
generation the whole trajectory is re-fitted by a single non-negative lasso
over batch-average prediction columns on a held-out fold; most batches receive
weight exactly zero and are physically deleted from the model. At the shipped
defaults this discards 77--92\% of the trees that were grown.

We give a complete description of the method and its implementation, and a
head-to-head evaluation against a \emph{tuned} LightGBM on nine datasets with a
common size-accounting rule applied to both sides. \SXE\ wins on accuracy on
four datasets and loses on five; it carries the smaller model on five of nine,
by up to $11.8\times$. Predictions are bit-identical across thread counts by
construction. The clearest result in the study is negative: \SXE's shipped
defaults give up a median 9.5\% error reduction to tuning where LightGBM's give
up 3.3\%, so the competitive column is a tuned \SXE. We analyse why the
accuracy wins and losses fall where they do --- the method is a
variance-reduction device, and it wins where variance dominates and loses where
bias does --- and we are explicit about what the evaluation does not establish.
\end{abstract}

\tableofcontents

\newpage

\section{Introduction}

\subsection{Two constructions and the gap between them}

Two ensemble constructions dominate supervised learning on tabular data, and
they differ in how their members are related to one another.

The \textbf{random forest} \cite{breiman2001rf} grows $M$ trees that are
mutually independent given the data. Each sees a bootstrap sample of the rows
and a random subset of features at each split, each is fitted to the
\emph{original} target, and the prediction is their unweighted average.
Independence is the point: averaging $M$ estimators with correlation $\rho$ and
variance $\sigma^2$ gives variance $\rho\sigma^2 + (1-\rho)\sigma^2/M$, so
decorrelating the members buys variance reduction that more members alone
cannot. The construction is embarrassingly parallel, and its bias is the bias
of a single deep tree --- which it does not attempt to reduce.

The \textbf{gradient-boosted decision tree} \cite{friedman2001gbm} grows trees
in strict sequence. Tree $m$ is fitted to the negative gradient of the loss at
the current ensemble prediction --- for squared error, the residual --- and is
added with a shrinkage factor $\nu$. Each member exists only as a correction to
its predecessors. This attacks bias directly and is why GBDTs dominate tabular
benchmarks \cite{grinsztajn2022}, but it is inherently sequential: tree $m{+}1$
cannot begin until tree $m$ has updated the residual. Modern implementations
\cite{ke2017lightgbm,chen2016xgboost} recover parallelism \emph{within} a tree
--- across features and histogram bins --- rather than across trees.

The gap between them is a design space, not a dichotomy. A boosting step need
not be one tree. If a step consists of $B$ trees grown independently against
the \emph{same} residual and then averaged, the step is a miniature random
forest, the residual update happens once per step rather than once per tree,
and the $B$ trees within a step can be grown concurrently. Nothing about
gradient boosting requires $B=1$; the choice is historical.

\subsection{What \SXE\ is}

\SXE\ (``Sparse eXtra Ensembles'') takes that observation and builds three
further ideas on it.

\begin{enumerate}[leftmargin=*]
\item \textbf{Batch boosting.} A boosting step grows $B$ randomized,
subsampled trees in parallel against the current residual and applies their
average as one shrunken correction. $B$ is a parameter (\code{batch\_size}),
and it interpolates continuously between stochastic gradient boosting ($B=1$)
and a random forest (one batch).

\item \textbf{Per-tree holdout pruning at randomized strength.} Each tree is
grown on a random fraction of the rows, so the complement is a holdout that
tree has never seen. \SXE\ routes it through the grown tree and runs classical
reduced-error pruning \cite{quinlan1987rep} bottom-up, collapsing any subtree
that does not beat its own collapsed leaf on holdout squared error by a
per-leaf penalty $\lambda$. Crucially $\lambda$ is drawn \emph{per tree} from a
log-uniform distribution, so the ensemble contains members ranging from
essentially unpruned to nearly stumps. No single bias--variance trade-off is
chosen at fit time; the final stage arbitrates.

\item \textbf{One-shot sparse non-negative re-weighting.} After generation,
\SXE\ discards the uniform weights that boosting assigned and re-fits the
trajectory with a non-negative lasso on a held-out fold, using one column per
\emph{batch} (the batch's average prediction) rather than one per tree. Most
coefficients come out exactly zero, and the corresponding batches are deleted
from the model object. This is Friedman and Popescu's ISLE construction
\cite{friedman2003isle} --- generate a dictionary of trees, then fit a sparse
linear combination of them --- with Breiman's non-negativity constraint from
stacked regressions \cite{breiman1996stacked}, applied at batch granularity.
\end{enumerate}

The per-tree randomization (depth, minimum leaf size, feature fraction,
subsample rate, split rule, pruning strength all drawn independently for every
member) is in the spirit of extremely randomized trees \cite{geurts2006et}: the
method averages over hyperparameter draws instead of committing to a single
setting. The histogram machinery --- quantile pre-binning, parent-minus-child
histogram subtraction, exclusive feature bundling, optional leaf-wise growth
--- follows LightGBM \cite{ke2017lightgbm}.

\subsection{What the design is for}

The method is not proposed as a general accuracy improvement over tuned GBDTs;
Section~\ref{sec:results} shows it is not one. It targets three properties that
follow from its structure:

\begin{description}[leftmargin=1.5em,style=nextline]
\item[Model size.] The lasso deletes most of the generated ensemble, and the
survivors were already pruned against honest holdout data. Where this lands
well the model is several times smaller than a tuned LightGBM's at comparable
accuracy.
\item[Reproducibility.] Every tree's random number stream is seeded from its
\emph{index} in the generation order, never from which worker thread picked it
up. The number of threads is therefore a pure performance knob: the same seed
produces bit-identical predictions on 1 core or 10.
\item[Parallel structure.] Trees within a batch are independent, so growth
parallelises across trees, which sequential GBDT cannot do. Only the
between-batch residual update is serial.
\end{description}

\subsection{Contributions}

\begin{enumerate}[leftmargin=*,itemsep=2pt]
\item A complete specification of the method (Section~\ref{sec:method}),
including the exact split criterion, the pruning inequality and its degenerate
case, the batch-boosting loop, and the weight fit.
\item An implementation account (Section~\ref{sec:impl}) covering the hybrid
histogram grower, the thread pool, the determinism argument, and the
serialization format.
\item A head-to-head evaluation against a tuned LightGBM on nine datasets under
one size-accounting rule, with paired per-seed analysis
(Section~\ref{sec:results}).
\item A quantified negative result: the shipped defaults are substantially
worse than the method (Section~\ref{sec:tuninggap}), with evidence about
\emph{which} defaults are wrong.
\item An analysis of where the method's accuracy wins and losses fall and why
(Section~\ref{sec:discussion}), and an explicit account of what the study does
not establish (Section~\ref{sec:limitations}).
\end{enumerate}

\section{Notation and setting}\label{sec:notation}

Training data are $(x_i, y_i)_{i=1}^{n}$ with $x_i \in \mathbb{R}^p$. Write
$X \in \mathbb{R}^{n\times p}$ for the design matrix and $y \in \mathbb{R}^n$
for the target.

The loss is squared error throughout. Regression minimises mean squared error
directly. Binary classification is handled as regression on $y \in \{0,1\}$ ---
the Brier-score reduction --- so leaf means are conditional-probability
estimates and a squared-error split on 0/1 targets is equivalent to a Gini
split up to constants. There is no logistic loss and there are no second-order
(Newton) leaf values; Section~\ref{sec:limitations} returns to what this costs.
Multiclass is a one-vs-rest reduction over the binary core.

\SXE\ partitions the training data once, before anything else happens. A
fraction $\rho_{\text{val}}$ (\code{validation\_fraction}, default $0.18$) is
set aside as a validation fold $(X_{\text{val}}, y_{\text{val}})$ of size
$n_{\text{val}}$; the remainder is the generation set on which all trees are
grown. The validation fold is used for early stopping, for the weight fit, and
for probability calibration.

The fitted model has the form
\begin{equation}\label{eq:model}
\hat{y}(x) \;=\; \beta_0 \;+\; \sum_{b \in \mathcal{S}} w_b \cdot
\underbrace{\frac{1}{B}\sum_{t \in \text{batch } b} T_t(x)}_{\text{batch mean}},
\qquad w_b \ge 0,
\end{equation}
where each $T_t$ is a pruned regression tree, $B$ is the batch size, and
$\mathcal{S}$ is the sparse set of batches surviving the weight fit. Members of
batches outside $\mathcal{S}$ are deleted from the model object rather than
merely zero-weighted, which is what makes the size claim about the artefact and
not about a hypothetical pruning.

\section{The method}\label{sec:method}

Figure-free, the pipeline is: bin once $\rightarrow$ (optionally bundle binary
columns) $\rightarrow$ generate a dictionary of trees by batch boosting, each
tree self-pruned $\rightarrow$ fit non-negative sparse weights on the
validation fold $\rightarrow$ delete everything with zero weight $\rightarrow$
prune the bin table to the columns that survive.

\subsection{Global quantile binning}\label{sec:binning}

Every feature is quantile-binned once, up front, and \emph{every} downstream
stage --- split search, holdout routing, pruning, member prediction, the
weight-fit design matrix --- operates on that single \code{uint8} matrix. There
is no intermediate float matrix.

\code{max\_bins} counts the bins available to finite values and is at most 255,
following scikit-learn's \code{HistGradientBoosting*} vocabulary; the 256th bin
is reserved for missing values and is not part of the budget. Edges are
estimated from a subsample of at most $50{,}000$ rows: the code takes
$\code{max\_bins}-1$ interior points of $\text{linspace}(0,1,\cdot)$ as quantile
levels, deduplicates the resulting edges (so a feature with few distinct values
gets few bins), and maps values with
\code{searchsorted(edges, x, side='right')}. The layout handed to the C++ core
is column-major, since the split search sweeps a feature at a time.

Binning is not merely a speed device here. A uniform random threshold drawn on
a binned feature is \emph{exactly} a random bin cut-point, so the
extremely-randomized split rule of Section~\ref{sec:growth} composes with
binning without approximation --- unlike a random threshold on raw values,
which binning would have to round.

\paragraph{Missing values.} \code{NaN} gets its own bin. At the default
\code{max\_bins}$=255$ the finite values occupy bins $0..254$ and missing
occupies 255, which needs no support in the tree code at all: bin 255 routes
like any other. Missing therefore travels with the largest values by default,
and a split at the top finite bin --- which the exact-cut scan finds whenever it
pays --- separates it. Infinity is still rejected at
\code{fit} and \code{predict}: \code{NaN} means missing, but infinity is a
number the binning cannot place. An all-missing column gets no edges at all and
every row lands in bin 255.

This is weaker than LightGBM's learned default direction, which can send missing
either way at every node, and the repository is explicit that it is a
convenience feature rather than an accuracy feature. The evidence is mixed by
design: on a synthetic problem where missingness \emph{is} the signal the scheme
reaches 0.97 AUC, which no imputation can reach; but against median imputation
at 10\% and 30\% missingness on four datasets it wins four of eight comparisons
and loses four, with Covertype the clearest loss --- punching holes in binary
indicator columns puts \code{NaN} above the positive class, where the median
puts it with the negatives. The point is that the library stops refusing the
data, not that it handles it better than the user would.

\subsection{Exclusive feature bundling}\label{sec:efb}

\code{BinaryBundler} implements LightGBM-style exclusive feature bundling
\cite{ke2017lightgbm} restricted to strictly binary columns, and is off by
default (\code{bundle\_binary\_features}). Columns whose values lie in
$\{0,1\}$ are ordered by descending nonzero count and greedily assigned to the
first bundle whose used-row mask they do not intersect; conflict tolerance is
zero and a bundle holds at most 250 members. A bundle of size $\ge 2$ becomes
one small-integer column: $0$ if no member is active, else the $1$-indexed
position of the active member. Singletons pass through unchanged.

The encoding is lossless up to ordinality. A threshold split on a bundle
selects a \emph{contiguous group} of the bundled categories rather than
isolating one, which is weaker than a one-vs-rest split, but trees compose
splits and the measured accuracy effect was nil. On Covertype this compresses
54 features to 12 --- bundles of 5 wilderness and 39 soil indicators --- and
cuts histogram work by roughly $4.5\times$.

Exclusivity is verified on the training data and must hold at predict time; a
prediction row that activates two members of one bundle would otherwise be
silently encoded as the first of them, indistinguishable from a row where only
that member is set. The current code raises, naming the offending row and
columns.

\subsection{Tree growth}\label{sec:growth}

Trees are grown by the C++ core on the binned matrix. For a node containing
$m$ rows with target sum $S$, a candidate split of feature $f$ at bin threshold
$t$ (routing ``bin $\le t$'' left) is scored by the variance-reduction proxy
\begin{equation}\label{eq:gain}
G(f,t) \;=\; \frac{S_L^2}{n_L} \;+\; \frac{S_R^2}{n_R} \;-\; \frac{S^2}{m},
\end{equation}
where $(n_L, S_L)$ and $(n_R, S_R)$ are the counts and target sums of the two
children. This is the standard reduction in sum of squares, written so that
only prefix sums of the node's histogram are needed. A candidate is invalid if
either child would fall below \code{min\_samples\_leaf}. A node stops splitting
when it reaches the tree's depth cap, when $m < 2\,\code{min\_samples\_leaf}$,
or when its target is numerically pure.

\paragraph{Candidate generation.} At each node, $K$ distinct candidate features
are drawn without replacement by a partial Fisher--Yates shuffle, with
$K = \text{round}(\phi p)$ and $\phi \sim U(\code{feat\_lo}, \code{feat\_hi})$
drawn \emph{once per tree}. For each candidate feature, cut-points are
evaluated in one of two modes:
\begin{description}[leftmargin=1.5em,style=nextline]
\item[Exact (\code{splitter="best"}, the default).] A prefix scan over the
feature's occupied bin range returns the maximiser of \eqref{eq:gain}.
\item[Extremely randomized (\code{splitter="random"}).] \code{n\_thresholds}
cut-points are drawn uniformly on the occupied bin range and only those are
scored.
\end{description}
With a retained histogram both cost $O(\text{bin span})$, so the exact scan is
essentially free; the choice is not about speed but about the trade between
split quality and member diversity, and it is genuinely dataset-dependent ---
the tuned configurations of Section~\ref{sec:tuninggap} chose random cuts on
the high-dimensional and one-hot datasets and exact cuts elsewhere.

\paragraph{Hybrid depth-wise growth with histogram subtraction.} Nodes with at
least \code{hybrid\_min} rows (default 600) retain a full all-feature histogram
of $p \times 256$ (count, sum) pairs. Candidates are then evaluated from prefix
sums with no pass over the data at all. After the node's index range is
partitioned in place, the \emph{smaller} child's histogram is built from data
in $O(|\text{small}| \cdot p)$ and the larger child's is obtained by an
in-place parent-minus-child subtraction, reusing the parent's buffer. At most
$\text{depth}+1$ buffers are live per thread and they are pooled.

Children below \code{hybrid\_min} fall back to a second grower that builds a
single 256-bin histogram per candidate feature with touched-bin clearing, and
so pays no $p\times256$ clearing cost on small nodes. The threshold trades one
cost against the other; 600 was found to work across all three of the datasets
it was tuned on and has not needed a per-dataset value since.

\begin{observation}[\code{hybrid\_min} is a pure performance knob]
Both growers select splits identically but draw from one shared per-tree random
stream, so the order in which they recurse determines which node a given random
draw lands on. Both therefore recurse \emph{smaller child first}. Without that,
changing \code{hybrid\_min} --- ostensibly a speed setting --- would silently
change the fitted model. The repository tests this invariance rather than
asserting it.
\end{observation}

\paragraph{Leaf-wise growth.} Setting \code{max\_leaf\_nodes} switches to
best-first growth in LightGBM's style: a max-heap keyed on gain holds pending
leaves, the highest-gain leaf is split, each child's best split is evaluated and
pushed if it clears the depth, size and gain floors, and growth stops at the
leaf budget. This is not a corner case --- the random search of
Section~\ref{sec:tuninggap} selected it on four of nine datasets.

A pending leaf carries only its best split, not the histogram that found it; the
histogram is rebuilt when the leaf is popped and released immediately, so at
most one is live per workspace. That gives up parent-minus-child subtraction for
the children of a popped node, costing roughly one extra histogram fill per
node, and it is nonetheless \emph{faster}: the alternative --- one retained
\code{Hist} per pending leaf, at $p \times 256 \times 12$ bytes each --- made the
workspace scale with the leaf budget rather than with depth, and its per-buffer
reset \code{memset} dominated the fill it was saving. On $4000 \times 500$ data
at \code{batch\_size}$=16$ the retained-histogram scheme reached 906\,MB peak RSS
in 0.90\,s against 254\,MB in 0.62\,s for the rebuild scheme, and 304\,MB for
depth-wise growth.

Exported node ids satisfy $\text{parent} < \text{children}$ in both modes, which
is what makes the bottom-up passes below valid without a separate traversal
order.

\paragraph{Per-tree randomization.} Every tree independently draws:
\begin{center}
\begin{tabular}{ll}
\toprule
Quantity & Distribution \\
\midrule
\code{max\_depth} & uniform over \code{max\_depth\_choices} ($0 \Rightarrow$ capped at 64) \\
\code{min\_samples\_leaf} & log-uniform on \code{min\_samples\_leaf\_range} \\
feature fraction $\phi$ & uniform on \code{max\_features\_range} \\
row subsample rate & uniform on \code{subsample\_range}, sampled without replacement \\
pruning strength $\lambda$ & log-uniform on \code{prune\_strength\_range} \\
\bottomrule
\end{tabular}
\end{center}
No two members of a batch share hyperparameters. This is the ExtraTrees idea
carried up a level: rather than choosing one configuration, the ensemble
averages over the draw.

\subsection{Holdout reduced-error pruning}\label{sec:pruning}

Each tree is grown on a subsample of rate drawn in $[0.4, 0.75]$ by default, so
25--60\% of the rows are \emph{not} used to grow it. These are that tree's
private holdout, and they are free --- the out-of-bag idea \cite{breiman1996oob}
applied to pruning rather than to scoring.

After growth, the holdout rows are routed through the tree once, accumulating
$(n_h, \sum y, \sum y^2)$ at \emph{every} node on each row's root-to-leaf path.
A single bottom-up pass --- descending node id, valid because ids are
pre-order --- then computes, for each node, the holdout sum of squared errors
that collapsing that node to a leaf would incur:
\begin{equation}\label{eq:sseleaf}
\SSE_{\text{leaf}} \;=\; \sum_h y^2 \;-\; 2\mu \sum_h y \;+\; n_h \mu^2,
\end{equation}
where $\mu$ is the node's \emph{training} mean --- the value the collapsed leaf
would actually predict. The node is collapsed if and only if
\begin{equation}\label{eq:prune}
\SSE_{\text{leaf}} \;\le\; \SSE_{\text{subtree}} \;+\;
\lambda\,\hat\sigma^2_h\,(L_{\text{subtree}} - 1),
\end{equation}
where $\SSE_{\text{subtree}}$ is the sum of the (possibly already collapsed)
children's values, $L_{\text{subtree}}$ is the subtree's post-collapse leaf
count, and $\hat\sigma^2_h$ is the holdout variance at the root, which makes
$\lambda$ scale-free.

The interpretation is direct. At $\lambda = 0$ the rule removes exactly those
splits that do not help on data the tree has never seen --- dead wood. As
$\lambda$ grows, each additional leaf must pay for itself at a rate of
$\lambda\hat\sigma^2_h$, giving a nested sequence of increasingly aggressive
subtrees: a holdout analogue of weakest-link cost-complexity pruning
\cite{breiman1984cart}, but with the complexity penalty evaluated against
honest data rather than resubstitution error.

Each tree emits exactly one member, pruned at its own
$\lambda \sim \text{log-}U(\lambda_{\text{lo}}, \lambda_{\text{hi}})$. Pruning
diversity therefore lives \emph{across} trees, and the weight fit of
Section~\ref{sec:weights} decides how much of each strength the model wants.
Collapsed nodes are exported as leaves carrying their training means.

\begin{observation}[The small-holdout degeneracy]\label{obs:guard}
If the subsample rate approaches 1, the holdout is nearly empty and many
subtrees are visited by no holdout row at all. For such a subtree both sides of
\eqref{eq:prune} are $0 \le 0 + \lambda\hat\sigma_h^2(L-1)$, which is satisfied
trivially, and the tree collapses --- in the worst case to a stump. \SXE\
guards this by emitting the tree unpruned whenever the holdout has fewer than
$\max(50, n/50)$ rows. The fix was worth $0.4504 \to 0.4437$ RMSE on California
housing in the original measurement, and it is what makes mild-subsampling
configurations viable at all.
\end{observation}

\subsection{Batch-boosted dictionary generation}\label{sec:boost}

The generation loop runs entirely in the C++ core. Let $r$ denote the residual
on the generation set and $F_{\text{val}}$ the running prediction on the
validation fold.

\begin{algorithm}[H]
\caption{Batch boosting (\code{Ensemble::boost})}
\begin{algorithmic}[1]
\State $r \gets y$; \quad $F_{\text{val}} \gets 0$; \quad
$\text{best} \gets \infty$; \quad $\text{bad} \gets 0$
\For{$b = 1, \dots, \code{max\_iter}$}
  \State \textbf{in parallel} for $t = 1, \dots, B$:
  \Statex \hspace{2em} seed tree $(b,t)$ from its index; draw its hyperparameters
  \Statex \hspace{2em} draw a row subsample; grow a tree on $(X, r)$
  \Statex \hspace{2em} route the complement; prune by \eqref{eq:prune} at its own $\lambda$
  \Statex \hspace{2em} evaluate $T_t$ on the generation rows and on $X_{\text{val}}$
  \State $r \gets r - \dfrac{\nu}{B}\displaystyle\sum_{t=1}^{B} T_t(X)$
  \Comment{one shrunken correction per batch}
  \State $F_{\text{val}} \gets F_{\text{val}} + \dfrac{\nu}{B}\displaystyle\sum_{t=1}^{B} T_t(X_{\text{val}})$
  \State $\text{mse} \gets \|y_{\text{val}} - F_{\text{val}}\|^2 / n_{\text{val}}$
  \If{$\text{mse} < \text{best} - \tau$} $\text{best} \gets \text{mse}$;\ $\text{bad} \gets 0$
  \ElsIf{$\texttt{++bad} \ge \code{n\_iter\_no\_change}$} \textbf{break}
  \EndIf
\EndFor
\end{algorithmic}
\end{algorithm}

The early-stopping tolerance is $\tau = \code{tol} \cdot \Var(y_{\text{val}})$,
so it is expressed in units of the target's variance rather than in absolute
error. The residual update is the only serial dependency in the loop; the $B$
tree constructions inside a batch touch no shared mutable state.

Two properties of the loop are worth naming. First, trees fit residuals, and
their holdout pruning is evaluated against the same residual target --- which
is the correct objective at that stage, since the tree's job is to explain
what remains. Second, the batch average $\frac{1}{B}\sum_t T_t$ is itself a
small random forest, so each boosting step is a shrunken correction built from
a variance-reduced estimate rather than from a single high-variance tree. This
is the sense in which $B$ dials between the two classical methods:

\begin{center}
\begin{tabular}{lll}
\toprule
\code{batch\_size} & Recovers & Character \\
\midrule
$1$ & stochastic gradient boosting & one tree per residual update \\
$10$--$16$ (default range) & \SXE & a mini-forest per residual update \\
$\to$ whole budget & random forest & one update, everything fitted to $y$ \\
\bottomrule
\end{tabular}
\end{center}

Empirically $B \in [4,16]$ preserved almost all of boosting's accuracy while
exposing the parallel width.

\subsection{The sparse non-negative weight fit}\label{sec:weights}

Generation produces up to $\code{max\_iter} \times B$ trees with uniform
implied weights $\nu/B$. \SXE\ discards those weights and re-fits.

Let $P \in \mathbb{R}^{n_{\text{val}} \times M}$ hold every member's prediction
on the validation fold. Columns are averaged within batches,
$P_b = \frac{1}{B}\sum_{t \in b} P_{\cdot t}$, giving one column per batch;
this averaging happens inside the C++ core so the design matrix is never
materialised at member granularity. The weights solve a positive elastic net
with intercept --- a pure lasso at the default \code{l1\_ratio}$=1$:
\begin{equation}\label{eq:lasso}
\min_{w \ge 0,\; \beta_0} \;
\frac{1}{2 n_{\text{val}}} \big\| y_{\text{val}} - \beta_0 - P_b w \big\|_2^2
\;+\; \alpha \|w\|_1 .
\end{equation}
The regularisation strength $\alpha$ is chosen by 3-fold cross-validation over
a path of 20 values, with Gram matrices precomputed --- which took the fit from
roughly 20\,s to 1--5\,s at thousands of columns and to 0.02--0.3\,s at
hundreds. Batches whose weight is zero are deleted; each surviving member
receives weight $w_b / B$, and $\beta_0$ becomes the model intercept. If the
cross-validated solution collapses entirely to the intercept, a tiny-$\alpha$
positive fit is substituted so the model remains an ensemble rather than a
constant.

Four design choices deserve justification.

\paragraph{Why a lasso at all.} By the stagewise--lasso correspondence
\cite{efron2004lars,hastie2007forward}, forward stagewise fitting with
infinitesimal steps traces the lasso path. The boosting trajectory is exactly
such a stagewise fit, so a lasso over the generated columns is a convex stand-in
for the greedy procedure that produced them --- it can express the same
solutions, and it can express sparser ones the greedy fit passed through and
never returned to. This is ISLE \cite{friedman2003isle}.

\paragraph{Why non-negative.} Breiman's stacked regressions
\cite{breiman1996stacked} found that constraining combination weights to be
non-negative yields sparser and markedly more stable solutions than
unconstrained least squares. It also prevents the pathology where two batches
acquire large opposing weights that cancel: every retained batch must
contribute in the direction it was fitted. Note that the weights are
\emph{not} additionally constrained to sum to one --- the fit is free to scale
the ensemble, and the intercept absorbs whatever shared bias remains.

\paragraph{Why on held-out data.} Selection is deliberately on the validation
fold. Fitted in-sample, the deepest and least-pruned members would absorb
everything and the pruning diversity of Section~\ref{sec:pruning} would be
wasted.

\paragraph{Why batch-average columns.} Cost. Per-member columns made the fit
explode to $\approx 7$\,s at thousands of columns; batch columns made it
$0.02$--$0.3$\,s with no measured accuracy change, and the fit's cost becomes
independent of the tree count. The price is granularity: a batch is retained or
dropped as a unit, and the model cannot express ``keep three of these ten
trees''.

\begin{observation}[The lasso compresses; it does not correct]
A diagnostic comparing the fitted weights against the raw uniform generation
weights $\nu/B$ found test-set-equal accuracy with roughly $10\times$
compression. The weight fit's role is therefore \emph{trajectory pruning} ---
deciding which steps of the boosting path are redundant --- rather than
repairing the direction of the fit. Generation decides where each batch points;
the lasso decides how much of it to keep, and whether to keep it at all.
\end{observation}

The prototype's earlier multi-round class (\code{CppSXE} in \code{reference/})
retains a different schedule for the same fit --- full cross-validation on the
first and last rounds, single warm fixed-$\alpha$ fits at $0.7\times$ the
previous $\alpha$ in between, and a narrowed final grid
$\alpha_1 \cdot \text{logspace}(0.5, -1, 10)$ --- which the one-shot fit
described here replaced.

Measured at the shipped defaults on two datasets (script in
Section~\ref{sec:repro}):

\begin{center}
\begin{tabular}{lrrrr}
\toprule
Dataset & Batches run & Trees grown & Trees kept & Kept \\
\midrule
California housing & 300 & 4{,}800 & 368 & 7.7\% \\
Friedman1 + noise & 155 & 2{,}480 & 560 & 22.6\% \\
\bottomrule
\end{tabular}
\end{center}

\noindent The California figure is not an accident of that dataset: the defaults
budget 300 batches and early stopping never fires, so the generator runs its full
course and the lasso then discards 92\% of what it built.

\subsection{Prediction, classification and calibration}\label{sec:predict}

Prediction routes binned rows through the retained trees and forms
$\beta_0 + \sum_j w_j T_j(x)$, parallelised over rows.

Binary classification predicts the same raw score. Because the score is a
weighted combination of 0/1 leaf means it is approximately a probability but
not a calibrated one, so \SXE\ fits a Platt sigmoid \cite{platt1999} on the
validation fold: $p = \sigma(a\,s + b)$ by maximum likelihood against Platt's
smoothed targets $(N_+{+}1)/(N_+{+}2)$ and $1/(N_-{+}2)$, which keep the
optimum finite when the fold happens to be separable. The smoothing is applied
by giving every observation both labels with weights $t$ and $1-t$, so an
ordinary logistic solver optimises the smoothed objective exactly. If the fold
cannot identify a calibration --- one class only, a constant score, or a fitted
map that is non-increasing --- the calibrator is dropped and the raw score is
clipped to $[10^{-6}, 1-10^{-6}]$ instead, so a bad fit can never reorder the
scores. Under the one-vs-rest reduction this decision is made for all classes
together: a rare class whose validation fold contains no positives cannot be
calibrated, and mapping the others while clipping that one would normalise a row
across two different scales.

Multiclass is one-vs-rest: one complete ensemble per class, with per-class
scores normalised to sum to one. This is a reduction, not a native objective,
and it costs a full ensemble per class --- on the 26-class Letter dataset, 26
ensembles.

\subsection{Determinism}\label{sec:determinism}

The random stream for the tree at batch $b$, slot $t$ is seeded as
\begin{equation}
\text{seed}_{b,t} \;=\; \texttt{splitmix64}\big(s \;+\; \texttt{0x9E3779B9} \cdot (bB + t + 1)\big),
\end{equation}
a function of the tree's \emph{index} only. Nothing about which worker thread
executes the tree, or in what order the pool drains its queue, enters the
stream. Combined with the smaller-child-first recursion of both growers, this
gives:

\begin{itemize}[leftmargin=1.5em,itemsep=1pt]
\item identical models across \code{n\_jobs} settings;
\item identical models across \code{hybrid\_min} settings;
\item identical models across repeated runs.
\end{itemize}

All three are covered by tests in the repository (\code{tests/test\_determinism.py}),
not merely asserted. The invariance is \emph{within} a platform. Per-tree
randomness is drawn through \code{std::uniform\_int\_distribution} and
\code{std::uniform\_real\_distribution}, whose algorithms are
implementation-defined by the C++ standard, so libstdc++, libc++ and MSVC
produce different sequences from the same seeded \code{mt19937\_64} and
therefore grow different trees. Making seeds portable would require replacing
the standard distributions with fixed implementations; the repository records
this as an open decision.

\subsection{Serialization and model size}\label{sec:size}

Two distinct quantities must be kept apart.

An earlier version of this comparison set an \SXE\ format estimate against
\emph{pickled} baselines, which overstated the difference by two to five times;
that claim was withdrawn, and the rule below is what replaced it.

\paragraph{The portable accounting formula.} For comparing tree structure, size
is computed identically for both sides: 13 bytes per internal node (feature 4,
threshold 1, left 4, right 4), 8 bytes per leaf (value), plus a full 255-edge
float64 bin table per input feature. Trees are full binary, so internal nodes
$=$ leaves $-$ trees. This is the column reported in Section~\ref{sec:results},
and it special-cases neither library --- but it compares \emph{structure} and
nothing else, and it is a poor guide to what either library would ship. It
charges LightGBM for a bin table its format does not need, and it charges both
sides for edges that deduplicate away on low-cardinality data or are never
indexed on wide data. The repository measures \SXE's real bin table at 1.9\,kB
on Letter where the formula charges 32.6\,kB, and at 67\,kB on Madelon where the
formula charges 1.02\,MB.

\paragraph{A second, less flattering column.} The benchmark script has since
been extended to record a \code{real\_bytes} column alongside the portable one:
measured for \SXE\ (\code{model\_size\_bytes\_}, the bytes \code{serialize()}
writes plus the quantile edges of the columns the trees actually index), and for
LightGBM the smallest sensible encoding of what its trees need --- a float64
threshold rather than a bin id, an explicit left child, and no bin table, at 20
bytes per internal node. That assumption is generous to LightGBM, whose real
files also carry split gains and node counts. It is nonetheless \emph{less}
favourable to \SXE\ than the portable formula: on Madelon LightGBM's real model
is the smaller one, which the portable column hid behind a bin table neither
library actually pays. The committed results predate this change, so
Section~\ref{sec:results} reports the portable column throughout and
Section~\ref{sec:limitations} carries the consequence.

\paragraph{The format \SXE\ actually writes.} \code{Ensemble::serialize} stores
only what prediction reads. Three fields that a naive layout would store are
recovered or discarded instead:
\begin{itemize}[leftmargin=1.5em,itemsep=1pt]
\item \emph{left child}: the exporter emits nodes in pre-order, pushing right
then left, so a node's left child is always the next id. It is recomputed on
load.
\item \emph{value on internal nodes}: prediction reads a value only at a leaf.
\item \emph{gain}: it existed for \code{feature\_importances\_}, which is now
accumulated once at fit time into a length-$p$ vector.
\end{itemize}
What remains is 9 bytes per internal node, 8 per leaf, and one bit per node
recording which is which. The bin table is additionally pruned to the columns
the retained trees actually index; on Madelon at defaults the model splits on
61 of 500 features and was carrying 421\,kB of quantile edges for the other
439. Unused columns transform to bin 0, which is safe precisely because nothing
reads them, and a one-vs-rest classifier shares a single table across its
sub-models covering the union of the columns they use, rather than each
sub-model holding a pruned copy of the same edges.

An earlier version of this format wrote 29 bytes per node --- storing
\code{gain} everywhere, \code{feat}/\code{thr}/\code{left}/\code{right} on
leaves and \code{val} on internal nodes --- so the advertised size was
$2.9\times$ smaller than the artefact. That has been corrected. Measured at the
shipped defaults:

\begin{center}
\begin{tabular}{lrrr}
\toprule
Dataset & \code{serialize()} & Portable formula & Ratio \\
\midrule
California housing & 472{,}757 B & 565{,}618 B & $0.84\times$ \\
Friedman1 + noise & 1{,}279{,}555 B & 1{,}542{,}667 B & $0.83\times$ \\
\bottomrule
\end{tabular}
\end{center}

\noindent \SXE's real tree format is therefore \emph{below} the portable
formula the benchmark charges it. Its bin-table charge runs the other way, and
by more, so the two errors do not cancel and neither sign can be claimed for the
size column as a whole.

\subsection{Complexity}\label{sec:complexity}

Let $n_s$ be a tree's subsample size, $n_h$ its holdout size, $d$ its depth,
and $M_b$ the number of batch columns.

\begin{center}
\begin{tabular}{ll}
\toprule
Stage & Cost \\
\midrule
Growth (with subtraction) & $O(n_s (p + d))$ per tree \\
In-place partitioning & $O(n_s d)$ per tree \\
Holdout routing & $O(n_h d)$ per tree \\
Pruning pass & $O(\text{nodes})$ per tree \\
Residual/validation update & $O((n + n_{\text{val}}) d)$ per tree \\
Weight fit & $O(n_{\text{val}} M_b^2)$ Gram, then coordinate descent on $M_b \times M_b$ \\
\bottomrule
\end{tabular}
\end{center}

The $O(n_s p)$ term in growth is dominated by the smaller-child histogram fills,
which sum to at most $n_s p / 2$ per level because the larger child is always
obtained by subtraction. Histogram memory is $O(\text{depth})$ per thread in
depth-wise growth --- at most $\text{depth}+1$ pooled buffers live --- and
$O(1)$, a single live buffer, in leaf-wise growth; neither scales with the leaf
budget or the tree count. Measured single-core grower throughput on Covertype
was $\approx 27.8$k leaves/s against LightGBM's $\approx 13.9$k, so wall-clock
differences between the libraries come from \emph{how much structure each
method builds}, not from per-leaf cost.

\section{Implementation}\label{sec:impl}

\subsection{Layout}

\begin{center}
\begin{tabular}{ll}
\toprule
Path & Role \\
\midrule
\code{src/sxe/\_core.cpp} & C++17 core: growers, pruning, boosting, prediction, serialization \\
\code{src/sxe/estimator.py} & \code{SXERegressor}, \code{SXEClassifier} \\
\code{src/sxe/binning.py} & \code{QuantileBinner} \\
\code{src/sxe/bundling.py} & \code{BinaryBundler} \\
\code{benchmarks/} & the head-to-head suite and its committed outputs \\
\code{reference/} & the original ctypes prototype, kept as a parity oracle \\
\bottomrule
\end{tabular}
\end{center}

The division of labour is deliberate: everything that runs once per fit and is
cheap (binning, bundling, the train/validation split, the lasso) stays in
Python where scikit-learn's implementations are already good; everything that
runs once per tree or once per row is in C++.

The prototype under \code{reference/} is not installed and is retained for two
purposes: as the source of the v1/v2 ablation figures in
Section~\ref{sec:results}, and as a parity oracle --- the packaged core
reproduces it bit-for-bit. It holds \code{sxe.cpp} (the original C API core with
OpenMP, including the v2 multi-round \code{sxe\_grow}), \code{cpp\_sxe.py} (a
ctypes wrapper exposing \code{CppSXE} for v2 and \code{CppSXEBoost} for v3),
\code{sparse\_extra.py} (the pure-Python v1 prototype built on scikit-learn's
\code{ExtraTreeRegressor} with cost-complexity-path variants), the v1/v2
benchmark scripts, and \code{results\_v3.csv}.

\subsection{Build}

The package builds with scikit-build-core and pybind11 and needs a C++17
compiler and CMake from source; at runtime it needs only NumPy ($\ge 1.22$) and
scikit-learn ($\ge 1.3$) on Python $\ge 3.9$, with \code{lightgbm} required only
by the benchmark scripts. Continuous integration builds and tests wheels for
five platforms across CPython 3.9--3.13.

Optimisation flags are deliberately portable: \code{-march=native} measured no
reproducible speedup on these datasets, which is consistent with the
memory-bound histogram loops of Section~\ref{sec:complexity}, so wheels give it
up for nothing.

\subsection{Threading}

Growth and prediction use a hand-rolled \code{std::thread} pool rather than
OpenMP, so there is no runtime to bundle --- notably no \code{libomp} step for
macOS users --- and \code{n\_jobs} is explicit. The pool holds
$\text{threads}-1$ workers and uses the calling thread as the last one. Work is
distributed by an atomic counter, so scheduling is dynamic; this is safe for
determinism precisely because each work item's random stream is index-derived
(Section~\ref{sec:determinism}). The first exception thrown by any worker is
captured, the counter is poisoned so remaining workers stop promptly, and the
exception is rethrown on the calling thread. The GIL is released for the whole
generation loop.

\code{n\_jobs=None} means \emph{all cores}, matching other gradient-boosting
libraries rather than scikit-learn's single-core default; negative values follow
the joblib convention.

\subsection{scikit-learn conformance}

\code{SXERegressor} and \code{SXEClassifier} pass
\code{sklearn.utils.estimator\_checks.check\_estimator} in full, so they compose
with \code{Pipeline}, \code{GridSearchCV}, \code{cross\_val\_score} and
\code{joblib} without adapters. Fitted models pickle and deep-copy --- the trees
travel through the serialization format of Section~\ref{sec:size} --- so they
survive the round trip into joblib workers. Parameter names follow
scikit-learn's \code{HistGradientBoosting*} vocabulary wherever the concepts
match. The repository's test suite (105 tests at the time of writing) covers
the API surface, determinism, persistence, multiclass, calibration, validation
and output quality.

\section{Experimental setup}\label{sec:setup}

\subsection{Datasets}

Nine datasets, chosen so that each stresses a different property --- including
several that stress properties \SXE\ does not have.

\begin{center}
\small
\setlength{\tabcolsep}{4pt}
\begin{tabular}{llll}
\toprule
Dataset & Shape & Task & What it tests \\
\midrule
California housing & $20{,}640 \times 8$ & regression & smooth low-dimensional signal, small $n$ \\
Friedman1 + 20 noise & $25{,}000 \times 25$ & regression & known signal buried in noise \\
Superconductivity & $21{,}263 \times 81$ & regression & many correlated informative features \\
Covertype 2-vs-rest & $30{,}000 \times 54$ & binary & the original reference point \\
Covertype 2-vs-rest, full & $581{,}012 \times 54$ & binary & large $n$, where GBDT scaling shows \\
Adult & $48{,}842 \times 14$ & binary & categoricals (cardinality 2--41), missing values \\
Bank marketing & $45{,}211 \times 16$ & binary & categoricals (cardinality 2--12) \\
Letter recognition & $20{,}000 \times 16$ & 26-class & the cost of the one-vs-rest reduction \\
Madelon & $2{,}600 \times 500$ & binary & $p \gg n$, mostly irrelevant features \\
\bottomrule
\end{tabular}
\end{center}

Five seeds each, except Covertype-full and Letter at three. The seed drives the
80/20 train/test split (stratified for classification), the model, and --- for
Friedman1 --- the data itself.

\subsection{Protocol}

\paragraph{Both libraries are tuned.} A win over a misconfigured baseline
establishes nothing, so both libraries are tuned per dataset by random search
over a fixed space with the same number of draws (10--24, larger where a fit is
cheap). Draw zero is always the library's own defaults, so a tuned row cannot
be worse than its default row on the selection fold. The search runs once per
dataset on seed 0 and the selected configuration is reused across seeds, so
neither library tunes against a seed it is scored on. The hand-tuned
configurations from the ablation history (Section~\ref{sec:results}) are entered
as additional warm starts. Selected configurations are committed to
\code{benchmarks/tuned\_configs.json}.

\paragraph{Common preprocessing.} Both libraries receive the same matrix:
categorical columns ordinal-coded, numeric missing values median-imputed from
the training split. This is necessary because \SXE\ at the reference commit
rejects non-finite input and has no categorical handling. A separate
\code{LightGBM-native} row runs LightGBM on raw \code{NaN}s with declared
categorical columns, which prices what the workaround costs.

\paragraph{One size rule, both sides.} Size is the portable formula of
Section~\ref{sec:size} applied to both libraries' structure counts.

\paragraph{Timing.} Wall-clock on one 10-core laptop; each library is reported
at its own best thread count. On this machine LightGBM's best is one thread on
seven of nine datasets and \SXE's is all cores on all nine.

\paragraph{One asymmetry, in LightGBM's favour.} LightGBM's early-stopping fold
doubles as its selection fold. \SXE\ early-stops on a fold carved from the
training half and is therefore selected on data it has not seen.

\subsection{Provenance}

Every figure in Section~\ref{sec:results} is recomputed directly from
\code{benchmarks/results\_suite.csv}, the committed raw output of
\code{benchmarks/bench\_suite.py}. Accuracy and structure are read at the lowest
thread count so that repeated thread settings do not double-count;
time is the mean at each library's best thread count.

\section{Results}\label{sec:results}

\subsection{Head-to-head, both tuned}

\begin{table}[h]
\centering
\small
\caption{Both libraries tuned. Mean $\pm$ standard deviation across seeds.
Bold marks the better mean.}
\setlength{\tabcolsep}{4pt}
\begin{tabular}{llccrrrr}
\toprule
& & \multicolumn{2}{c}{Accuracy} & \multicolumn{2}{c}{Size (MB)} & \multicolumn{2}{c}{Best fit (s)} \\
\cmidrule(lr){3-4}\cmidrule(lr){5-6}\cmidrule(lr){7-8}
Dataset & Metric & LightGBM & \SXE & LGBM & \SXE & LGBM & \SXE \\
\midrule
California      & rmse $\downarrow$ & \textbf{0.4409} $\pm$ 0.0095 & 0.4480 $\pm$ 0.0097 & 2.845 & \textbf{0.241} & 2.22 & \textbf{1.40} \\
Friedman1+noise & rmse $\downarrow$ & 1.0695 $\pm$ 0.0112 & \textbf{1.0557} $\pm$ 0.0166 & \textbf{0.526} & 0.993 & \textbf{1.50} & 2.00 \\
Superconduct    & rmse $\downarrow$ & \textbf{9.3313} $\pm$ 0.3171 & 9.6685 $\pm$ 0.2807 & 5.766 & \textbf{1.253} & 13.33 & \textbf{9.63} \\
Covertype 30k   & auc $\uparrow$    & 0.9547 $\pm$ 0.0031 & \textbf{0.9590} $\pm$ 0.0026 & \textbf{1.092} & 3.545 & \textbf{1.32} & 8.32 \\
Covertype full  & auc $\uparrow$    & \textbf{0.9977} $\pm$ 0.0001 & 0.9946 $\pm$ 0.0001 & \textbf{10.508} & 14.448 & \textbf{50.31} & 82.11 \\
Adult           & auc $\uparrow$    & \textbf{0.9305} $\pm$ 0.0016 & 0.9240 $\pm$ 0.0020 & 0.489 & \textbf{0.443} & 1.72 & \textbf{1.69} \\
Bank marketing  & auc $\uparrow$    & 0.9338 $\pm$ 0.0017 & \textbf{0.9354} $\pm$ 0.0021 & 0.202 & \textbf{0.181} & \textbf{1.02} & 2.24 \\
Letter (26)     & acc $\uparrow$    & \textbf{0.9661} $\pm$ 0.0018 & 0.9629 $\pm$ 0.0003 & \textbf{5.109} & 13.564 & 67.61 & \textbf{33.85} \\
Madelon         & auc $\uparrow$    & 0.8873 $\pm$ 0.0089 & \textbf{0.9031} $\pm$ 0.0075 & 1.364 & \textbf{1.036} & \textbf{3.66} & 5.96 \\
\bottomrule
\end{tabular}
\end{table}

\noindent
\textbf{Accuracy: four to \SXE, five to LightGBM. Size: five to \SXE, four to
LightGBM. Time: three to \SXE, six to LightGBM.} No axis is a systematic win.

The timing column deserves one note about what it is not. Reported this way ---
each library at its own best thread count --- speed goes to \SXE\ on three of
nine. At the setting a user is actually likely to choose, all cores for both,
\SXE\ would win most of them, because LightGBM's within-tree parallelism costs
it on data this size. That framing is available and is declined here: a library
compared against itself at its worst configuration is the same error as a tuned
model compared against an untuned baseline.

\subsection{Paired analysis}

Means and standard deviations across seeds understate what a paired comparison
can say, because both libraries see the same split at each seed. Table
\ref{tab:paired} reports the per-seed difference.

\begin{table}[h]
\centering
\small
\caption{Paired per-seed differences, \SXE\ minus LightGBM. $t$ is the mean
difference in units of its own standard deviation; positive $t$ favours \SXE\
for AUC and accuracy, negative $t$ favours \SXE\ for RMSE.}
\label{tab:paired}
\begin{tabular}{llrrrl}
\toprule
Dataset & Metric & Mean diff & sd & $|t|$ & Verdict \\
\midrule
Covertype full  & auc  & $-0.00313$ & 0.00005 & 57.8 & LightGBM, decisive (0/3 seeds) \\
Adult           & auc  & $-0.00645$ & 0.00066 & 9.8  & LightGBM, clear (0/5) \\
Covertype 30k   & auc  & $+0.00436$ & 0.00088 & 5.0  & \SXE, clear (5/5) \\
Superconduct    & rmse & $+0.33720$ & 0.10956 & 3.1  & LightGBM, clear (0/5) \\
Madelon         & auc  & $+0.01574$ & 0.00661 & 2.4  & \SXE, clear (5/5) \\
Letter          & acc  & $-0.00317$ & 0.00166 & 1.9  & LightGBM, suggestive (0/3) \\
Bank marketing  & auc  & $+0.00157$ & 0.00090 & 1.7  & \SXE, suggestive (5/5) \\
California      & rmse & $+0.00712$ & 0.00629 & 1.1  & within noise (0/5) \\
Friedman1+noise & rmse & $-0.01380$ & 0.01265 & 1.1  & within noise (4/5) \\
\bottomrule
\end{tabular}
\end{table}

Two \SXE\ wins (Covertype-30k, Madelon) and three LightGBM wins (Covertype-full,
Adult, Superconductivity) are clear. Bank marketing and Letter are suggestive
but not decisive. California and Friedman1 are within seed-to-seed spread,
though the sign is consistent across seeds in each case. Note that the sign of
the mean difference is consistent with the win count on every dataset, so no
result is an artefact of one outlying seed.

\subsection{Model size}

Size is smaller for \SXE\ on five of nine, and where it wins it wins large:
$11.8\times$ on California, $4.6\times$ on Superconductivity. But the
mechanism deserves stating plainly, because it is not ``\SXE\ builds smaller
models''. Neither random search penalises size. Both columns therefore report
where an accuracy-driven search happened to land, and LightGBM's search drifts
to very large models --- 135k leaves on California, 267k on Superconductivity,
497k on Covertype-full --- whenever that buys it anything at all. Where \SXE\
loses on size it loses on structure it genuinely spends: 164k leaves on
Covertype-30k, and 650k on Letter, where the one-vs-rest reduction pays for 26
separate ensembles.

One caution on reading these numbers by proxy: leaf count is not a format-free
stand-in for size. The two libraries put different amounts of information in a
node and a leaf, and \SXE\ additionally carries a bin table that LightGBM does
not, so a leaf-count ratio and a byte ratio can point in different directions.

The five--four split is also specific to the portable accounting. Under the
real-format column of Section~\ref{sec:size} the Madelon verdict reverses,
because the portable formula charges both libraries for a 500-feature bin table
that \SXE\ prunes to 61 columns and LightGBM never needs at all. Read the size
column as a structure comparison, which is what it is.

\subsection{What tuning buys each library}\label{sec:tuninggap}

The gap between a library's tuned and default rows is the measurable form of
``does this need tuning''. Table~\ref{tab:tuning} reports the error reduction
from default to tuned, where error is RMSE for regression and $1 - \text{metric}$
for classification.

\begin{table}[h]
\centering
\small
\caption{Error reduction from each library's defaults to its tuned
configuration. Larger means the defaults were further from what the library can
do.}
\label{tab:tuning}
\begin{tabular}{lrrr}
\toprule
Dataset & \SXE & LightGBM & Ratio \\
\midrule
California      & 4.4\%  & 1.0\%  & $4.3\times$ \\
Friedman1+noise & 1.9\%  & 3.3\%  & $0.6\times$ \\
Superconduct    & 11.8\% & 2.4\%  & $5.0\times$ \\
Covertype 30k   & 50.4\% & 8.2\%  & $6.2\times$ \\
Covertype full  & 82.8\% & 61.4\% & $1.3\times$ \\
Adult           & 4.5\%  & 2.7\%  & $1.7\times$ \\
Bank marketing  & 2.6\%  & 0.9\%  & $2.9\times$ \\
Letter          & 36.8\% & 6.7\%  & $5.5\times$ \\
Madelon         & 9.5\%  & 3.6\%  & $2.6\times$ \\
\midrule
\textbf{median} & \textbf{9.5\%} & \textbf{3.3\%} & $\mathbf{2.9\times}$ \\
\bottomrule
\end{tabular}
\end{table}

\textbf{\SXE\ gains more from tuning on eight of nine datasets.} This is the
clearest result in the study and it is negative. It is a statement about the
shipped defaults rather than about the method --- the tuned column of
Table~\ref{tab:paired} is competitive, the default one is not --- but a user's
first fit is what they actually get, so the distinction is cold comfort.

The winning configurations say specifically what is wrong. Not one of the nine
used the default \code{max\_depth\_choices} $=(4,5,6,7,8)$: five chose
$(8,10,12)$ or $(10,12,14)$ and four chose leaf-wise growth with 63--127
leaves. \code{max\_iter} selected the top of its search range on seven of nine
and \code{batch\_size} on six. Three of the knobs that matter most therefore
ran into the ceiling of the search space, which means the defaults are not
merely mistuned but \emph{systematically build too small a model}, and the true
optimum lies beyond what was searched.

Those defaults have since been revised in exactly that direction
(Section~\ref{sec:inflight}), recovering a median 58\% of the per-dataset
tuning gain on six datasets re-checked, at roughly $4\times$ the fit time and
2--3.5$\times$ the model size. The table above is measured at the previous
settings and is left as the diagnosis that motivated the change; what the
revision is worth across the whole suite is a question for the next full run.

\subsection{What the missing features cost}

\begin{center}
\begin{tabular}{lrrr}
\toprule
Dataset & LightGBM (shared matrix) & LightGBM-native & Difference \\
\midrule
Adult          & 0.9305 auc & 0.9304 auc & $-0.0001$ \\
Bank marketing & 0.9338 auc & 0.9351 auc & $+0.0013$ \\
\bottomrule
\end{tabular}
\end{center}

\paragraph{RandomForest context.} For scale, the v1-era measurements against
scikit-learn's \code{RandomForest} (single seed, prototype, not re-run):
California 0.511 RMSE in 33.8\,s, Friedman1 1.443 RMSE in 116\,s, Covertype
0.9486 AUC in 11.5\,s. Every version of \SXE\ from v1 onward beats those on
both accuracy and time by wide margins, which is the floor the method clears
rather than a competitive comparison. The v1-era model sizes were pickled
scikit-learn objects and are not comparable to the size column above.

\medskip
Native categorical handling and native missing-value handling buy LightGBM
almost nothing on these two datasets. \SXE's gap on Adult ($-0.0065$ AUC) is
therefore a \emph{model} gap, not an encoding gap --- a useful negative result,
since it says that adding native categorical support would not close it. What
the missing features cost is user effort: the preprocessing must happen, and
today the user writes it.

\subsection{Ablations}

The following deltas were measured on the original prototype, under
Linux/libstdc++, single-seed. They are retained as the \emph{relative} evidence
for each design decision --- the direction and rough magnitude of each change
--- and will not reproduce exactly against Section~\ref{sec:results}, which is
a multi-seed re-measurement on macOS/libc++.

\begin{center}
\begin{tabular}{p{0.30\textwidth}p{0.63\textwidth}}
\toprule
Change & Effect \\
\midrule
Multi-round boosting vs.\ a single-round dictionary
& The single-round dictionary lands at random-forest level; the residual rounds
supply essentially all accuracy beyond bagging. Boosted multi-round beat RF
everywhere (Friedman 1.088 vs.\ 1.443 RMSE) at 20--40$\times$ smaller models;
LightGBM matched it only on the noisy task. \\
\addlinespace
Holdout pruning replacing a cost-complexity approximation
& Covertype AUC $0.9480 \to 0.9528$, alongside 5--10$\times$ faster growth from
the move to C++, and a $6\times$ cut in the weight fit from precomputing the
Gram matrix. Consistent with pruning decisions made against genuinely held-out
data. \\
\addlinespace
Histogram subtraction
& $\approx 1.6\times$ on growth. \\
\addlinespace
Exclusive feature bundling
& A further $\approx 2.4\times$ on Covertype at unchanged accuracy. Profiling at
this stage showed the grower already producing $\approx 2\times$ more leaves/s
than LightGBM, so the remaining wall-clock gap was structure volume and pipeline
overhead rather than per-leaf cost. Covertype became a clean three-axis win;
Friedman was accuracy-tied but slow, because the dictionary was a fixed size
where LightGBM early-stopped; California plateaued near 0.47 RMSE. \\
\addlinespace
Batch boosting + batch-group lasso (v2 $\to$ v3)
& Friedman flipped to a win on all three axes --- a 0.9-second grower and a
0.12\,MB model. Per-member lasso columns cost
$\approx 7$\,s at thousands of columns; batch columns made the fit
0.02--0.3\,s with no accuracy change. The shallow-tree defaults broke Covertype
outright (AUC 0.92), recovered only by depth 10--14 with random cuts: that
dataset genuinely needs deep, wide-candidate trees, and it was the first clear
sign of the defaults problem of Section~\ref{sec:tuninggap}. \\
\addlinespace
Small-holdout guard
& Observation~\ref{obs:guard}. California $0.4504 \to 0.4437$ RMSE. \\
\addlinespace
Leaf-wise growth
& Did not break the California accuracy plateau (63 leaves $\approx$ depth-wise
best, 31 leaves worse) --- evidence that the residual gap there is not about
tree shape. \\
\bottomrule
\end{tabular}
\end{center}

\section{Discussion}\label{sec:discussion}

\subsection{Why the wins and losses fall where they do}

Every device in \SXE\ is a variance-reduction device. Averaging $B$ independent
trees per step reduces the variance of each correction. Randomizing depth, leaf
size, feature fraction and subsample rate decorrelates the members.
Reduced-error pruning against a private holdout removes structure that does not
generalise. The non-negative lasso deletes whole steps of the trajectory that
do not pay on held-out data. There is no corresponding device for bias: the
loss is squared error, leaf values are training means, and there is no Newton
step.

The results follow that structure closely. \SXE\ wins where variance dominates
the error:
\begin{itemize}[leftmargin=1.5em,itemsep=1pt]
\item \textbf{Madelon} ($p \gg n$, 2{,}600 rows, 500 features, mostly
irrelevant) --- the regime where a single greedy learner overfits the
irrelevant features and averaging over feature subsets pays most. \SXE\ wins by
2.4 sd on 5/5 seeds, with a model $1.3\times$ smaller.
\item \textbf{Covertype-30k} --- a subsampled problem where the same argument
applies at moderate scale. 5.0 sd, 5/5 seeds.
\item \textbf{Friedman1 with 20 pure-noise features} --- signal buried in noise;
\SXE\ leads on 4/5 seeds, though within spread.
\end{itemize}

And it loses where bias dominates:
\begin{itemize}[leftmargin=1.5em,itemsep=1pt]
\item \textbf{Covertype-full} (465k training rows) is the sharpest case. With
that much data there is enough signal to fit the bias away, and a sequential
leaf-wise learner does that better than an averaged forest of randomized trees.
The gap is small in absolute terms ($0.0031$ AUC) but 57.8 sd --- utterly
consistent.
\item \textbf{Superconductivity} (81 correlated informative features) ---
randomizing the feature subset hurts when most features carry signal, because
the diversity it buys costs split quality that the data would have supported.
\item \textbf{California housing} (smooth, low-dimensional, strong signal) sits
at a plateau that neither leaf-wise growth nor deeper trees moved, which points
at leaf \emph{values} rather than tree \emph{shape} as the binding constraint
--- consistent with squared-error training means being the remaining
approximation.
\end{itemize}

This is a coherent story rather than a set of scattered outcomes, and it is
actionable: the levers most likely to close the remaining gaps are all on the
bias side (Section~\ref{sec:future}).

\subsection{The lasso's real role}

It is tempting to read the weight fit as a corrective stage that improves on
the boosting trajectory. The diagnostic says otherwise: fitted weights and raw
uniform generation weights give equal test accuracy, at roughly $10\times$
compression. The lasso is a \emph{compression} stage. That is exactly what the
method needs it to be --- the size claim is its entire justification --- but it
also means the accuracy of \SXE\ is the accuracy of the batch-boosted
generator, and improvements should be sought there.

\subsection{An observation on fold reuse}

The single validation fold does four jobs: early stopping, the lasso weight
fit, the cross-validation of $\alpha$ inside that fit, and Platt calibration.
The $\alpha$ selection is properly nested within the fold, so that is sound.
But the batches retained by the lasso are selected on the same fold that early
stopping already optimised against, so the retained set --- and therefore the
model size, which is the headline claim --- is chosen on data the generator has
already been steered toward. The test sets are untouched, so
Section~\ref{sec:results}'s accuracy figures are not affected; the concern is
specifically that the compression ratio is measured under mild optimism. A
second fold, or a cross-fitted weight fit, would settle it.

\section{Limitations and threats to validity}\label{sec:limitations}

Stated at the strength the evidence supports.

\paragraph{The defaults are not the method.} Section~\ref{sec:tuninggap} is the
dominant caveat. Everything in the \SXE\ column of Table~\ref{tab:paired} is a
tuned \SXE, and at the settings those results were measured under, a user's
first fit gave up a median 9.5\% error reduction. The defaults have since been
revised and recover a median 58\% of that on the datasets re-checked
(Section~\ref{sec:inflight}), but the revision has not been measured across the
whole suite, so the honest description of the package remains ``competitive
after a search, and better out of the box than these tables show''.

\paragraph{Size is an outcome, not a controlled variable.} Both searches
optimise accuracy alone, so a single tuned point per library cannot answer the
question the use case actually poses --- \emph{what accuracy is available if the
model has to fit in $N$ megabytes}. The search has since been instrumented to
record every draw's (loss, size) pair, which answers it at no extra cost, but
the results reported here predate that and report only where an accuracy-driven
search happened to land.

\paragraph{The size column compares structure, not artefacts.} The portable
accounting is applied identically to both sides, which is the right call for
comparing tree structure, but it is a byte count of neither library. It charges
LightGBM for a bin table its format does not need, and it charges both sides for
quantile edges that deduplicate away or are never indexed --- \SXE's real bin
table is 1.9\,kB on Letter against the formula's 32.6\,kB. The net effect is not
uniformly in \SXE's favour: its own tree format is measurably \emph{below} the
formula (Section~\ref{sec:size}), but its bin-table charge is far above what it
ships, and under the real-format accounting the Madelon size verdict reverses.
The correct reading is that structure is compared fairly and shipped size is
not yet compared at all.

\paragraph{Nine datasets, one machine.} Timings are wall-clock on a 10-core
laptop and will not transfer. Three seeds on the two most expensive datasets is
thin for a variance estimate.

\paragraph{Seeds are not portable across platforms.} Bit-identical
reproducibility holds across runs, thread counts and \code{hybrid\_min} on one
platform, and is tested. It does not hold across C++ standard libraries
(Section~\ref{sec:determinism}).

\paragraph{Method limits.} Squared error throughout; no logistic loss; no
second-order leaf values. Multiclass is a one-vs-rest reduction costing one
ensemble per class --- on Letter, 26 ensembles, 650k leaves and $2.7\times$
LightGBM's model for slightly worse accuracy. Features are capped at 255 finite
bins plus the missing bin. There is no general categorical handling beyond
binary bundling, and no \code{sample\_weight}.

\paragraph{Preprocessing was done for the benchmark.} Categoricals were
ordinal-coded and missing values imputed before either library saw them.
Section~\ref{sec:results} prices what that costs in accuracy --- almost nothing
on this suite --- but not what it costs in user effort, which is real. \SXE\ now
takes \code{NaN} directly (Section~\ref{sec:binning}), so of that burden only
the categorical encoding remains; the benchmark predates the change and imputed
for both libraries anyway, which keeps the comparison model-against-model.

\paragraph{Concurrency.} Concurrent \code{fit} calls on a single estimator
instance are unsupported, as they are for scikit-learn estimators generally.

\section{Related work}\label{sec:related}

\SXE\ is a recombination of established components rather than a new primitive,
and its lineage is worth stating precisely.

The \textbf{dictionary-then-sparse-combination} structure is Friedman and
Popescu's Importance Sampled Learning Ensembles \cite{friedman2003isle}:
generate a large family of trees by some randomized procedure, then fit a
regularised linear combination of them. RuleFit \cite{friedman2008rulefit} is
the same idea at the granularity of individual rules (root-to-leaf paths)
rather than whole trees. \SXE\ sits at a \emph{coarser} granularity than either
--- one column per batch --- for cost reasons, and Section~\ref{sec:future}
notes the finer end as unexplored.

The \textbf{non-negativity constraint} on combination weights is Breiman's
stacked regressions \cite{breiman1996stacked}, motivated there by exactly the
stability and sparsity properties \SXE\ relies on.

The \textbf{lasso-as-stagewise} justification is the LARS correspondence
\cite{efron2004lars} and its forward-stagewise analysis
\cite{hastie2007forward}: the boosting path and the lasso path coincide in the
infinitesimal-step limit, which is what licenses fitting a convex problem in
place of the greedy one.

The \textbf{pruning} is classical reduced-error pruning \cite{quinlan1987rep}
against honest holdout data, with a complexity penalty in the style of
cost-complexity pruning \cite{breiman1984cart} but evaluated out of sample. The
use of the out-of-subsample complement is the out-of-bag device
\cite{breiman1996oob} applied to pruning rather than to error estimation.

The \textbf{per-tree randomization} is in the spirit of extremely randomized
trees \cite{geurts2006et}, extended from the split threshold to the tree's
whole hyperparameter vector.

The \textbf{histogram machinery} --- pre-binning, parent-minus-child histogram
subtraction, exclusive feature bundling, leaf-wise growth --- is LightGBM's
\cite{ke2017lightgbm}, as is much of the parameter vocabulary via
scikit-learn's \code{HistGradientBoosting*}.

What is not, to our knowledge, standard is the specific combination: batch
boosting as a continuous interpolation between forests and boosting;
\emph{randomizing the pruning strength across members} so the ensemble spans a
bias--variance spectrum that a later convex fit arbitrates; and applying the
sparse non-negative fit at batch granularity so that it prunes the boosting
\emph{trajectory} rather than individual members.

\section{Future work}\label{sec:future}

\subsection{Unbuilt levers already paid for}

Three refinements target leaf \emph{values} rather than tree structure, which
Section~\ref{sec:discussion} identifies as where the remaining gap most
plausibly lives. All three are cheap given what the pipeline already computes.

\begin{description}[leftmargin=1.5em,style=nextline]
\item[Honest leaf values.] The holdout is already routed through every tree and
its per-node $(n_h, \sum y, \sum y^2)$ are already accumulated for the pruning
decision --- but exported leaves carry \emph{training} means. Re-estimating
leaf values from the holdout, as honest trees do in the causal-forest
literature \cite{wager2018causal}, is nearly free given what is in hand.
\item[Hierarchical shrinkage.] Shrink each node's prediction toward its parent
in proportion to sample counts, then delete nodes whose post-shrinkage
contribution vanishes \cite{agarwal2022hs}. Reported to often beat hard pruning
on accuracy, and it composes with the existing bottom-up pass rather than
replacing it.
\item[Leaf-level re-weighting.] The weight fit operates on batch-average
columns, coarser than per-tree. The fine end is a ridge or lasso over the
leaf-membership matrix --- RuleFit territory \cite{friedman2008rulefit} ---
which can express corrections that per-batch weights cannot.
\end{description}

\subsection{Then, in expected order of value}

\begin{enumerate}[leftmargin=1.5em,itemsep=1pt]
\item \textbf{Fix the defaults.} Section~\ref{sec:tuninggap} names the
direction: deeper trees, more batches, larger batches. Because three knobs hit
the ceiling of the search space, the first step is a wider search, not a
re-pick within the current one. This is the highest-value work in the list by a
distance, because it is the difference between the method and the product.
\item \textbf{Newton (second-order) leaf values and a learning-rate schedule},
both aimed at the bias-side gap.
\item \textbf{Logistic loss and a native multiclass objective}, which would
retire the one-vs-rest reduction and the $26\times$ cost on Letter.
\item \textbf{A size-budgeted evaluation}, so the size claim answers the
question the use case poses.
\item \textbf{A GPU port}, for which the batch structure is unusually
well-suited: within a batch, all trees $\times$ all nodes at a depth $\times$
all candidates reduce to gather and histogram kernels with no sort or argmax
over raw values, and the weight fit is one small GEMM plus coordinate descent.
\end{enumerate}

\subsection{Changes landed since the benchmark run}\label{sec:inflight}

Four changes have landed since the results of Section~\ref{sec:results} were
measured. None of them alters the method; two of them change what a user gets,
and the committed benchmark output predates all four.

\emph{Missing-value support}, as described in Section~\ref{sec:binning}. This
removes the preprocessing burden priced in Section~\ref{sec:results} ---
though, as that section shows, the preprocessing was costing accuracy almost
nothing, so what it removes is user effort.

\emph{A defaults revision.} \code{learning\_rate} $0.15 \to 0.07$,
\code{max\_iter} $120 \to 300$, \code{batch\_size} $10 \to 16$,
\code{n\_iter\_no\_change} $8 \to 15$, and \code{max\_depth\_choices}
$(4,5,6,7,8) \to (6,8,10,12)$ --- one configuration derived from the pattern in
Section~\ref{sec:tuninggap} rather than tuned per dataset, moving all three of
the identified knobs the way the tuned configurations pointed. On six datasets
re-checked it improves every one and recovers a median 58\% of the per-dataset
tuning gain. It is not free: fits take roughly $4\times$ longer and default
models are 2--3.5$\times$ larger, which trades directly against the size claim.
That is the right trade while default accuracy is the larger gap, and the five
parameters can be set back. A full re-run of the suite is what would settle it;
per the project's own rule, a default change is a product decision measured
across the whole benchmark, not on one dataset.

\emph{The dual size accounting} of Section~\ref{sec:size}, together with a
record of every tuning draw's (loss, size) pair in \code{tuning\_draws.csv}, so
the size-budget question of Section~\ref{sec:limitations} becomes answerable on
the next full run.

\emph{Two benchmark-protocol changes.} Covertype-full is now timed at all cores
only, as Letter already was --- a single-threaded fit at the new defaults on
465k rows runs to ten minutes, and the thread sweep would dominate the suite for
a comparison nobody would make at that size --- and the \SXE\ search cap rises
from 120 to 200 batches so sampled draws are not held below the default
trajectory length.

\emph{A leaf-wise memory fix} (Section~\ref{sec:growth}) and a publishing path:
a \code{v*} tag now pushes wheels to PyPI by trusted publishing, though the
package is not yet installable from PyPI.

\section{Reproduction}\label{sec:repro}

\begin{lstlisting}[language=bash]
pip install -e ".[bench]"

python benchmarks/bench_suite.py                  # full suite, ~1.5 h
python benchmarks/bench_suite.py -d adult letter  # a subset
python benchmarks/bench_suite.py --seeds 0        # single seed, quick
python benchmarks/bench_suite.py --retune         # redo both random searches

python -m pytest tests/                           # 105 tests
\end{lstlisting}

Raw output lands in \code{benchmarks/}\allowbreak\code{results\_suite.csv} and
selected configurations in
\code{benchmarks/}\allowbreak\code{tuned\_configs.json}; both are committed as
the provenance for every table in Section~\ref{sec:results}. Datasets download
via scikit-learn; Friedman1 is synthetic with a fixed seed. Figures reproduce
exactly only on the same C++ standard library
(Section~\ref{sec:determinism}).

The compression and serialized-size measurements in
Sections~\ref{sec:weights} and~\ref{sec:size} were taken at the shipped defaults
with:

\begin{lstlisting}[language=Python]
m = SXERegressor(random_state=0).fit(X, y)
nodes, leaves = m.ensemble_.structure(m._members)
blob     = len(m.ensemble_.serialize())   # == m.tree_bytes_
portable = (nodes - leaves) * 13 + leaves * 8
\end{lstlisting}

\noindent The prototype benchmarks in \code{reference/} still run from a
checkout, and are the source of the v1/v2 figures in
Section~\ref{sec:results}:

\begin{lstlisting}[language=bash]
cd reference
g++ -O3 -march=native -fopenmp -shared -fPIC sxe.cpp -o libsxe.so
python benchmark.py        # v1: RandomForest + LightGBM baselines + Python SXE
python cpp_benchmark.py    # v2: multi-round C++ SXE
python parity_check.py     # the packaged core against the prototype
\end{lstlisting}

\section{Conclusion}

\SXE\ is a coherent point in the design space between the random forest and the
gradient-boosted machine, reached by asking what happens when a boosting step
is a small forest rather than a single tree. The three devices built on that
--- randomized members, per-member holdout pruning at randomized strength, and
a non-negative lasso that deletes most of the trajectory --- are all
variance-reduction devices, and the empirical results track that fact closely:
the method wins where variance dominates and loses where bias does, on a
head-to-head against a tuned LightGBM that splits four--five on accuracy and
five--four on size.

The engineering claims are stronger than the accuracy claim and are the more
durable part of the contribution. Reproducibility across thread counts is
structural rather than best-effort, because per-tree seeds derive from tree
indices and both growers recurse in the same order; it is tested, not asserted.
The serialized artefact is now smaller than the portable formula used to
compare libraries, so the advertised size is conservative. And the parallel
structure --- independent trees within a batch, one serial update between
batches --- is a genuine architectural advantage over sequential GBDT that also
maps unusually cleanly onto a GPU.

The honest headline is the negative one. The shipped defaults give up a median
9.5\% error reduction to tuning where LightGBM's give up 3.3\%, and the tuned
configurations agree on the direction of the fix: deeper trees, more batches,
larger batches. The method is competitive; the product is not yet. Closing that
specific gap is worth more than any of the accuracy refinements on the list,
because a user's first fit is what they actually get.

\begin{thebibliography}{99}
\small

\bibitem{agarwal2022hs}
A.~Agarwal, Y.~S. Tan, O.~Ronen, C.~Singh, and B.~Yu.
\newblock Hierarchical shrinkage: Improving the accuracy and interpretability
of tree-based methods.
\newblock In \emph{ICML}, 2022.

\bibitem{breiman1996oob}
L.~Breiman.
\newblock Out-of-bag estimation.
\newblock Technical report, Statistics Department, UC Berkeley, 1996.

\bibitem{breiman1996stacked}
L.~Breiman.
\newblock Stacked regressions.
\newblock \emph{Machine Learning}, 24(1):49--64, 1996.

\bibitem{breiman2001rf}
L.~Breiman.
\newblock Random forests.
\newblock \emph{Machine Learning}, 45(1):5--32, 2001.

\bibitem{breiman1984cart}
L.~Breiman, J.~Friedman, R.~Olshen, and C.~Stone.
\newblock \emph{Classification and Regression Trees}.
\newblock Wadsworth, 1984.

\bibitem{chen2016xgboost}
T.~Chen and C.~Guestrin.
\newblock XGBoost: A scalable tree boosting system.
\newblock In \emph{KDD}, 2016.

\bibitem{efron2004lars}
B.~Efron, T.~Hastie, I.~Johnstone, and R.~Tibshirani.
\newblock Least angle regression.
\newblock \emph{Annals of Statistics}, 32(2):407--499, 2004.

\bibitem{friedman2001gbm}
J.~H. Friedman.
\newblock Greedy function approximation: A gradient boosting machine.
\newblock \emph{Annals of Statistics}, 29(5):1189--1232, 2001.

\bibitem{friedman2003isle}
J.~H. Friedman and B.~E. Popescu.
\newblock Importance sampled learning ensembles.
\newblock Technical report, Stanford University, 2003.

\bibitem{friedman2008rulefit}
J.~H. Friedman and B.~E. Popescu.
\newblock Predictive learning via rule ensembles.
\newblock \emph{Annals of Applied Statistics}, 2(3):916--954, 2008.

\bibitem{geurts2006et}
P.~Geurts, D.~Ernst, and L.~Wehenkel.
\newblock Extremely randomized trees.
\newblock \emph{Machine Learning}, 63(1):3--42, 2006.

\bibitem{grinsztajn2022}
L.~Grinsztajn, E.~Oyallon, and G.~Varoquaux.
\newblock Why do tree-based models still outperform deep learning on typical
tabular data?
\newblock In \emph{NeurIPS Datasets and Benchmarks}, 2022.

\bibitem{hastie2007forward}
T.~Hastie, J.~Taylor, R.~Tibshirani, and G.~Walther.
\newblock Forward stagewise regression and the monotone lasso.
\newblock \emph{Electronic Journal of Statistics}, 1:1--29, 2007.

\bibitem{ke2017lightgbm}
G.~Ke, Q.~Meng, T.~Finley, T.~Wang, W.~Chen, W.~Ma, Q.~Ye, and T.-Y. Liu.
\newblock LightGBM: A highly efficient gradient boosting decision tree.
\newblock In \emph{NeurIPS}, 2017.

\bibitem{platt1999}
J.~Platt.
\newblock Probabilistic outputs for support vector machines and comparisons to
regularized likelihood methods.
\newblock In \emph{Advances in Large Margin Classifiers}, 1999.

\bibitem{quinlan1987rep}
J.~R. Quinlan.
\newblock Simplifying decision trees.
\newblock \emph{International Journal of Man-Machine Studies}, 27(3):221--234,
1987.

\bibitem{wager2018causal}
S.~Wager and S.~Athey.
\newblock Estimation and inference of heterogeneous treatment effects using
random forests.
\newblock \emph{Journal of the American Statistical Association},
113(523):1228--1242, 2018.

\end{thebibliography}

\appendix

\section{Parameter reference}\label{app:params}

Shipped defaults. Parameters follow scikit-learn's
\code{HistGradientBoosting*} vocabulary where the concepts match; the third
column gives the name used by the prototype in \code{reference/}, for reading
that code. The five values marked $\dagger$ were revised after the benchmark run
(Section~\ref{sec:inflight}); the results of Section~\ref{sec:results} were
measured at the previous settings, shown in brackets.

\begin{center}
\footnotesize
\setlength{\tabcolsep}{3pt}
\begin{tabular}{lllp{0.36\textwidth}}
\toprule
Parameter & Default & Prototype & Meaning \\
\midrule
\multicolumn{4}{l}{\emph{Boosting trajectory}} \\
\code{learning\_rate} & \code{0.07}$^\dagger$ \small[0.15] & \code{shrink} & $\nu$, shrinkage on each batch-average correction \\
\code{max\_iter} & \code{300}$^\dagger$ \small[120] & \code{max\_batches} & maximum boosting batches \\
\code{batch\_size} & \code{16}$^\dagger$ \small[10] & \code{batch} & trees grown in parallel per batch \\
\code{n\_iter\_no\_change} & \code{15}$^\dagger$ \small[8] & \code{patience} & early-stopping patience, in batches \\
\code{tol} & \code{1e-4} & --- & improvement threshold, in units of $\Var(y_{\text{val}})$ \\
\code{validation\_fraction} & \code{0.18} & \code{val\_fraction} & fold for early stopping, weights and calibration \\
\addlinespace
\multicolumn{4}{l}{\emph{Per-tree randomization}} \\
\code{max\_depth\_choices} & \code{(6,8,10,12)}$^\dagger$ & \code{depth\_choices} & depth sampled uniformly; \code{0} $\Rightarrow$ 64 \\
& \small[(4,5,6,7,8)] & & \\
\code{min\_samples\_leaf\_range} & \code{(10, 60)} & \code{leaf\_range} & log-uniform per tree \\
\code{max\_features\_range} & \code{(0.5, 1.0)} & \code{feat\_range} & candidate-feature fraction, uniform per tree \\
\code{subsample\_range} & \code{(0.4, 0.75)} & --- & row subsample; complement is the pruning holdout \\
\code{prune\_strength\_range} & \code{(1e-3, 1.0)} & \code{lam\_range} & $\lambda$, log-uniform per tree \\
\addlinespace
\multicolumn{4}{l}{\emph{Tree construction}} \\
\code{splitter} & \code{"best"} & \code{argmax\_bins} & exact bin scan, or \code{"random"} cut-points \\
\code{n\_thresholds} & \code{2} & \code{n\_thresh} & random cut-points when \code{splitter="random"} \\
\code{max\_leaf\_nodes} & \code{None} & \code{max\_leaves} & if set, leaf-wise growth with this budget \\
\code{max\_bins} & \code{255} & \code{n\_bins} & bins for finite values; a 256th holds missing \\
\code{bundle\_binary\_features} & \code{False} & \code{bundle\_binary} & exclusive feature bundling for binary columns \\
\code{hybrid\_min} & \code{600} & --- & node size above which subtraction is used (speed only) \\
\addlinespace
\multicolumn{4}{l}{\emph{Weight fit and output}} \\
\code{l1\_ratio} & \code{1.0} & \code{l1\_ratios} & \code{1.0} is a pure lasso \\
\code{n\_alphas} & \code{20} & \code{n\_alphas} & regularisation path length \\
\code{weight\_cv} & \code{3} & \code{cv} & folds for choosing $\alpha$ \\
\code{calibrate} & \code{"sigmoid"} & --- & Platt scaling; \code{None} clips to $[10^{-6}, 1-10^{-6}]$ \\
\addlinespace
\multicolumn{4}{l}{\emph{Execution}} \\
\code{n\_jobs} & \code{None} & --- & threads; \code{None} $\Rightarrow$ all cores. Does not affect the model \\
\code{random\_state} & \code{None} & --- & seeds binning, the split, and per-tree randomization \\
\bottomrule
\end{tabular}
\end{center}

\section{Tuning guidance}\label{app:tuning}

Distilled from the nine per-dataset searches. Across all of them the same three
moves recurred --- \textbf{deeper trees, more batches, a larger batch size} ---
and no winning configuration used the default depth choices.

\begin{description}[leftmargin=1.5em,style=nextline]
\item[Smooth, low-dimensional signal.] Fine trajectories: small
\code{learning\_rate}, many batches, small \code{batch\_size}, mild subsampling
(0.85--1.0), \code{splitter="best"}, \code{max\_leaf\_nodes} $\approx 63$.
\item[High-dimensional or one-hot data.] \code{max\_depth\_choices = (10,
12, 14)}, \code{splitter = "random"} with \code{n\_thresholds = 3},
\code{bundle\_binary\_features = True}, moderate subsampling.
\item[Noisy data with irrelevant features.] Aggressive subsampling (0.2--0.4)
and a large \code{learning\_rate} with tight \code{n\_iter\_no\_change}. This is
also where the speed comes from.
\end{description}

\section{Diagnostics exposed by a fitted model}

\begin{lstlisting}[language=Python]
model.summary()             # batches, trees grown/kept, leaves, size, timings
model.feature_importances_  # split-gain, weighted by ensemble weight
model.n_iter_               # batches actually run before early stopping
model.n_trees_, model.n_trees_kept_
model.model_size_bytes_     # the bytes serialize() writes, not an estimate
\end{lstlisting}

With \code{bundle\_binary\_features=True}, a bundle's gain is divided equally
among its member columns --- a threshold on a bundle separates a group of them,
not one --- so those entries are a per-bundle figure spread over its columns
rather than a per-column measurement.

\end{document}
